%%=============================================================================
%% Toolselectie en vergelijkende evaluatie
%%=============================================================================

\chapter{\IfLanguageName{dutch}{Toolselectie en vergelijkende evaluatie}{Tool selection and comparative evaluation}}%
\label{ch:fase2}

\section{Scope van de tool-evaluatie}

Voor dit onderzoek stonden twee automatische netwerkdetectie-aanpakken binnen de praktische haalbaarheid centraal:

\begin{itemize}
  \item \textbf{NetBox Discovery} als primaire discovery-oplossing, aangevuld met \textbf{Diode} om discoverydata in NetBox te synchroniseren binnen de context van Citymesh.
  \item \textbf{Slurp'it} (via een NetBox plugin), dat belovend leek voor het verzamelen en structureren van netwerkdata, maar dat in de context van dit bachelorproject financieel minder haalbaar bleek doordat de oplossing betalend is.
\end{itemize}

Andere discovery-tools werden in deze casestudy niet volledig uitgevoerd of uitgerold tot op productieniveau. Dit betekent dat de vergelijking in dit hoofdstuk vooral bestaat uit: (1) een concrete evaluatie van de NetBox Discovery-workflow (met Diode) en (2) een meer kwalitatieve beoordeling van Slurp'it op basis van wat getest en geobserveerd kon worden.

De \texttt{NetworkTool}-applicatie (met NetBox API) wordt in dit hoofdstuk beschouwd als een \emph{aanvulling} op de discoveryflow (om ontbrekende datatypes zoals LLDP-neighbors/kabelrelaties te capteren), maar niet als een zelfstandige discovery-tool die in dezelfde zin als NetBox Discovery en Slurp'it vergeleken werd.

\section{Evaluatiecriteria}

Om te beoordelen of een oplossing voldoet aan de behoeften van Citymesh, werden volgende criteria in beschouwing genomen:

\begin{itemize}
  \item \textbf{Dekking:} kan de oplossing netwerkapparaten en relevante configuratiedetails bijhouden in een heterogene, dynamische omgeving?
  \item \textbf{Integratie met NetBox:} ondersteunt de oplossing een workflow die de \emph{single source of truth} in NetBox effectief versterkt?
  \item \textbf{Beheerlast:} vergt de oplossing veel configuratie/onderhoud of blijven herhaalde handmatige correcties beperkt?
  \item \textbf{Wireless relevantie:} kan de oplossing zinvol omgaan met draadloze use-cases (o.a. access points en devices in een distributed context)?
  \item \textbf{Topologie-informatie:} ondersteunt de oplossing, direct of via uitbreidbaarheid, relaties zoals interfaces en (indien mogelijk) kabel-/neighbor-relaties.
  \item \textbf{Kosten/haalbaarheid:} past de oplossing binnen een praktisch budget en is ze operationeel haalbaar voor een bachelorproef-casestudy.
  \item \textbf{Uitbreidbaarheid:} kan de oplossing geconfigureerd of aangevuld worden wanneer bepaalde datatypes (bv. neighbors/kabels) niet volledig automatisch landen?
 \end{itemize}

\section{Vergelijkende evaluatie}

\subsection{NetBox Discovery met Diode}

NetBox Discovery in combinatie met Diode werd als primaire oplossing beschouwd omdat dit binnen de scope van het onderzoek (gratis/low-cost) realiseerbaar was en direct inzetbaar is in NetBox.

In de praktijk bleek NetBox Discovery effectief voor het opbouwen van een bruikbare inventarisbasis: devices en hun documentatiedata konden via de discovery-workflow in NetBox terechtkomen. Voor het volledige SSOT-verhaal bleef echter één beperking over: LLDP-neighbors (\emph{neighbors}) werden niet automatisch als relaties in NetBox aangemaakt via de Orb/NetBox discoveryketen. Dit is vervolgens opgevangen met een aanvullende automatisering (\texttt{NetworkTool} + NetBox API) in hoofdstuk~\ref{ch:fase4}.

\subsection{Slurp'it (NetBox plugin)}

Slurp'it werd meegenomen omdat het in studies en documentatie veelbelovend oogt voor het ophalen en structureren van netwerkdata en omdat de beschikbare NetBox plugin een route biedt om discoverydata te integreren.

Tijdens de evaluatie werd vastgesteld dat Slurp'it conceptueel kan aansluiten op het doel (automatisch/gestructureerd netwerkdocumentatie), maar dat de oplossing betalend is en daardoor minder haalbaar wordt in de context van dit project. Hierdoor bleef Slurp'it buiten scope voor een volwaardige implementatie en werd het niet geselecteerd als primaire tool.

Slurp'it hanteert een licentiemodel waarbij het aantal ondersteunde devices per jaar bepalend is. De gratis versie ondersteunt slechts ongeveer 10 devices, terwijl Citymesh in de praktijk met veel hogere aantallen toestellen werkt (orde van grootte ~25.000 devices).

Voor een lage tot middelmatige schaal blijkt Slurp'it al snel duur: volgens de beschikbare prijsinformatie kost een licentie ongeveer 15.000 dollar per jaar voor minder dan 1000 devices. Voor meer dan 1000 devices moet bovendien een offerte aangevraagd worden, wat de praktische haalbaarheid in deze casestudy verder verlaagt.

Daar tegenover staat dat Slurp'it veel templates en een grote knowledge base bevat, waardoor met weinig inspanning toch goede resultaten kunnen bekomen worden voor het correct structureren en exporteren van netwerkdata. In combinatie met de device-licentielimieten was Slurp'it echter niet realistisch om als primaire oplossing te implementeren.

\section{Keuze voor de primaire oplossing}

Op basis van de scope, evaluatiecriteria en de praktische haalbaarheid werd gekozen voor \textbf{NetBox Discovery} (met Diode) als primaire oplossing. De belangrijkste redenen zijn:

\begin{itemize}
  \item \textbf{Haalbaarheid binnen budget:} NetBox Discovery was gratis en direct beschikbaar binnen de casestudy-scope.
  \item \textbf{Effectieve integratiebasis voor SSOT:} de workflow levert een bruikbare NetBox-inventaris op, waardoor de aanvullende automatisering kan starten van een “gevulde” basis.
  \item \textbf{Uitbreidbaarheid:} waar discovery niet volledig automatisch landt (LLDP-neighbors), kon je de ontbrekende functionaliteit aanvullen met een eigen pipeline (\texttt{NetworkTool} + NetBox API).
\end{itemize}

\section{Beperkingen van de vergelijking}

Doordat Slurp'it niet als primaire tool uitgerold werd en omdat andere tools niet in dezelfde mate als NetBox Discovery werden onderzocht of getest in de Citymesh-context, is dit hoofdstuk vooral een evaluatie van \emph{realistische opties} binnen de casestudy. De resultaten zijn daardoor minder “breed” dan een volledige marktstudie, maar wel concreet en toepasbaar in de eigen context van Citymesh.

